\documentclass[10pt,a4paper]{amsart}
\usepackage[margin=19mm]{geometry}
\usepackage{amsmath,amssymb,mathtools}
\usepackage[scr=rsfs,cal=euler]{mathalfa}
\usepackage{xfrac}
\usepackage[unicode,hidelinks,bookmarks=false]{hyperref}
\newcommand{\ie}{i.e.\ }
\newcommand{\cf}{cf.\ }
\newcommand{\resp}{respectively\ }
\newcommand{\Subsection}{\subsection}
\providecommand{\nobreakdash}{\nobreak\hbox{-}}
\setlength{\emergencystretch}{2em}
\multlinegap0pt
\usepackage{bm}
\usepackage{bbm}
\usepackage{dsfont}
\usepackage{xspace}
\usepackage{cancel}
\usepackage{mathtools}
\usepackage{amsthm}
\makeatletter
\@ifundefined{theorem}{\newtheorem{theorem}{Theorem}}{}
\@ifundefined{definition}{\newtheorem{definition}[theorem]{Definition}}{}
\makeatother
\newcommand\ad{\operatorname{ad}}
\newcommand\ip{{\langle\cdot {,} \cdot\rangle}}
\newcommand\n{{\mathfrak{n}}}
\newcommand\s{{\mathfrak{s}}}
\newcommand{\Ric}{\operatorname{Ric}}
\newcommand{\fraka}{\ensuremath{\mathfrak{a}} }
\newcommand{\frakn}{\ensuremath{\mathfrak{n}} }
\newcommand{\fraks}{\ensuremath{\mathfrak{s}} }
\newcommand{\smallfrac}[2]{\textstyle{\frac{#1}{#2}}  }
\renewcommand\a{{\mathfrak{a}}}
\title[Attached Submanifolds Beyond Symmetric Spaces]{Attached Submanifolds Beyond Symmetric Spaces}
\author{Megan M. Kerr, Tracy L. Payne}
\date{Source: arXiv:2601.13461v1 (CC BY 4.0)}
\begin{document}
\maketitle

\noindent\emph{Source and licence.} Attached Submanifolds Beyond Symmetric Spaces. Version arXiv:2601.13461v1 (CC BY 4.0), distributed under the Creative Commons Attribution 4.0 International licence, as recorded in the arXiv licence metadata for this exact version and shown on the abstract page https://arxiv.org/abs/2601.13461v1. Full rights notice: licenses/arxiv-2601.13461-CC-BY-4.0.txt.\par\smallskip

\noindent\textit{Editorial scope.} Counted unit: the Theorem whose source label is \texttt{thm: ricci-n-diff} in the article (source0.tex lines 728--741, source label \texttt{thm: ricci-n-diff}), delivered together with the article's own complete original proof of it (source0.tex lines 743--814, one \texttt{proof} environment of 72 lines). No mathematics is written, repaired or re-derived: statement and proof are the article's own text line for line. Required context retained uncounted: eq:jacobi-star (lines 322--327); eqn: ricci-nilpotent (lines 462--462); defn: attached (lines 546--553). Every definition, display and label of those lines is retained unchanged. Omitted from this package and not redistributed: 1--321, 328--461, 463--545, 554--727, 742--742, 815--1359. Nothing inside a retained excerpt is omitted or altered: every retained body is delivered exactly as the article's lines are written, so no omission receipt of any kind accompanies this package. Citations: the retained excerpts carry no citation command at all, so no bibliography entry of the article is transcribed and no reference is redistributed. Reference adaptation: none was needed and none was made; every reference inside the delivered unit resolves inside it, so no unresolved reference or citation is produced. Printed slips kept literal: no printed word, symbol, hypothesis or number of the counted statement or of its proof was altered. Layout and class: the article's own class is \texttt{amsart} with packages including amscd, amssymb, amsthm, amsxtra, array, mathtools, pdfsync, url, orcidlink; the wrapper is the lane's pinned \texttt{amsart} editorial wrapper at 10pt on A4, which restates the theorem-like environments the retained text needs and adds the article's own macro definitions that the retained text needs (\texttt{\textbackslash Ric}, \texttt{\textbackslash a}, \texttt{\textbackslash ad}, \texttt{\textbackslash fraka}, \texttt{\textbackslash frakn}, \texttt{\textbackslash fraks}, \texttt{\textbackslash ip}, \texttt{\textbackslash n}, \texttt{\textbackslash s}, \texttt{\textbackslash smallfrac}), with extra wrapper packages (bm, bbm, dsfont, xspace, cancel, mathtools). The delivered numbering is the unit's own. Assets: no retained span contains a figure environment, a graphics-inclusion command, a PDF-page inclusion, a table or a TikZ picture; no image file is copied into this package, the package declares no asset dependency and no image file is redistributed with it. Licence: version arXiv:2601.13461v1 of this article is distributed under the Creative Commons Attribution 4.0 International licence, as recorded in the arXiv licence metadata for this exact version and shown on the abstract page https://arxiv.org/abs/2601.13461v1. A full rights notice is delivered at licenses/arxiv-2601.13461-CC-BY-4.0.txt. Characters: every retained excerpt is pure ASCII, so the delivered engine, XeLaTeX with its default Latin Modern text font, reports no missing character.
\begin{equation}\label{eq:jacobi-star}
\smallfrac{1}{2} \sum _j \epsilon_j  [(\ad_{E_j^\perp})^{\ast, \frakn} ,
\ad_{E_j^\perp}
]
(X) =
\ad_{ \sum  \epsilon_j (\ad_{E_j^\perp})^{\ast, \fraks} E_j^\perp} (X) \end{equation}

 \begin{equation}\label{eqn: ricci-nilpotent}\Ric^{\n} = \tfrac 14 \sum \epsilon_i (\ad_{E_i}) \circ (\ad_{E_i})^* - \tfrac 12 \sum \epsilon_i (\ad_{E_i})^* \circ (\ad_{E_i}).\end{equation}

\begin{definition}\label{defn: attached} 
For a solvable pseudo-Riemannian metric  Lie algebra of strong Iwasawa 
type $(\s, \ip)$ and $\Lambda'$ as above,  the metric Lie algebra 
$(\s', \ip')$ is  the {\em attached subalgebra defined by $\Lambda'$.}  
Let $(S,g)$ be the simply connected  solvmanifold  associated to 
$(\s, \ip)$.  The submanifold $(S',g')$ defined by $(\s', \ip')$ is 
the {\em attached submanifold defined by $\Lambda'$. } 
\end{definition}

\begin{theorem}\label{thm: ricci-n-diff} 
   Let $(\fraks = \a \oplus \n,\ip)$ be a solvable  pseudo-Riemannian metric Lie algebra
   of strong Iwasawa type. 
   Let $(\fraks'= \fraka' \oplus \frakn', \ip')$ be an attached metric subalgebra defined by a set of roots $\Lambda'$
   as in Definition \ref{defn: attached}.
   Let $\{E_j'\} \cup \{E^\perp_j\}$ be a compatible orthonormal basis for $\n = \n' \oplus \n_0$.
   Then for $X$ in $\n',$  
   \[ \Ric^\n (X)  - \Ric^{\n'} (X)   = {\smallfrac{1}{2}} 
 \sum  \epsilon_j \left(  (\ad_{E_j^\perp}) \circ  (\ad_{E_j^\perp})^* -   (\ad_{E_j^\perp})^* \circ  (\ad_{E_j^\perp}) \right) (X).  
\]
 If in addition the Jacobi Star Condition holds, then
\[ \Ric^{\n} (X) - \Ric^{\n'} (X)  =  \ad_{H-H'}(X). 
\]
\end{theorem}

\begin{proof} We  compute  $\Ric^\n (X) - \Ric^{\n'} (X)$ for $X$ in $\n'.$
When we compute  $\Ric^{\n'},$ we will use that the  map $(\ad_{E_i^\prime})^{\ast, \frakn'}:\frakn' \to \frakn'$ 
 with respect to the restricted scalar product on $\frakn'$ is
equal to $\pi_{\n'} \circ (\ad_{E'_i})^{\ast,\frakn}$, where the adjoint map is relative to $(\frakn,\ip).$
By Equation \eqref{eqn: ricci-nilpotent},
\begin{align} \Ric^{\n} (X) &= \tfrac 14 \sum \epsilon_i  (\ad_{E_i}) \circ (\ad_{E_i})^* (X) - \tfrac 12 \sum \epsilon_i (\ad_{E_i})^* \circ  (\ad_{E_i})  (X), \, \text{and}
\notag \\
 \Ric^{\n'} (X) &= \tfrac 14 \sum \epsilon_i  (\ad_{E'_i}) \circ
 \pi_{\frakn^\prime} ((\ad_{E'_i})^* (X)) 
 - \tfrac 12 \sum \epsilon_i \, \pi_{\n'}((\ad_{E'_i})^* \circ  (\ad_{E'_i})  (X)). \notag 
\end{align}
We rewrite $ \Ric^{\n} (X)$ to sum over $\{ E_i^\prime\}$ and $\{E_i^\perp\}$ separately, to get
\begin{align*}  \Ric^{\n} (X) &= \tfrac 14 \left(\sum \epsilon_i  (\ad_{E_i^\perp}) \circ (\ad_{E_i^\perp})^* (X) + 
\sum \epsilon_i  (\ad_{E'_i})  \circ (\ad_{E'_i})^* (X) \right)    \\
& \qquad - \tfrac 12\left(\sum\epsilon_i  (\ad_{E_i^\perp})^* \circ  (\ad_{E_i^\perp})  (X) + 
\sum \epsilon_i  (\ad_{E'_i})^* \circ  (\ad_{E'_i})  (X)  \right). 
\end{align*}
We further expand  the second sum in each line above, using the decomposition
$(\ad_{E_i^\prime})^\ast(X) = \pi_{\frakn^\prime}
( (\ad_{E_i^\prime})^\ast (X) ) + 
\pi_{\frakn_0} ( (\ad_{E_i^\prime})^\ast(X) )$
to obtain
\begin{align}  \Ric^{\n} (X) 
&= \tfrac 14\left(A + \sum \epsilon_i  (\ad_{E'_i})  \circ \pi_{\n'}((\ad_{E'_i})^*
  (X)) +    B \right) 
\notag \\
& \qquad - \tfrac 12 \left(C + \sum  \epsilon_i \, \pi_{\n'}((\ad_{E'_i})^* \circ
  (\ad_{E'_i}) (X))  + D \right),\notag \end{align}
where $A, B, C$ and $D$ are defined by 
  \begin{alignat}{2} 
A &= \sum \epsilon_j (\ad_{E_j^\perp}) \circ  (\ad_{E_j^\perp})^* (X),& \qquad  
B &=  \sum \epsilon_i  (\ad_{E'_i}) \circ \pi_{\n_0}((\ad_{E'_i})^* (X)), \notag \\
C &=  \sum\epsilon_j  (\ad_{E_j^\perp})^* \circ  (\ad_{E_j^\perp})  (X),& \qquad
D &= \sum \epsilon_i \, \pi_{\n_0}((\ad_{E'_i})^* \circ (\ad_{E'_i}) (X)). \notag
\end{alignat}
The difference  $\Ric^{\n} (X) - \Ric^{\n'} (X) $ may  be written as 
\[ \Ric^{\n} (X) - \Ric^{\n'} (X)  = \smallfrac 14 A + \smallfrac 14  B - \smallfrac 12 C -\smallfrac 12 D. \]
We now simplify the terms $A, B, C$ and  $D$.  
First observe  
\begin{equation}\label{little identity} \langle (\ad_{E'_i})^*X, E_j^\perp \rangle = \langle X, [E'_i,
E_j^\perp] \rangle = -\langle X, [E_j^\perp , E'_i] \rangle = -\langle
(\ad_{E_j^\perp})^* X, E'_i \rangle.\end{equation}  
To begin, we show $B = A$: 
\begin{alignat*}{2}
 B &= \sum_i \epsilon_i  (\ad_{E'_i}) \circ \pi_{\n_0} ((\ad_{E'_i})^* X) \qquad &
  \text{by definition}\\ 
& = 
\sum_i\epsilon_i  [ E'_i, \sum_j \epsilon_j \langle (\ad_{E'_i})^* X, E_j^\perp \rangle
E_j^\perp] \qquad &  \strut \text{using the basis for $\frakn_0$}  \\
&=  \sum_i \epsilon_i  \sum_j \epsilon_j [ \langle (\ad_{E_j^\perp})^* X, E'_i \rangle
E_j^\perp, E'_i ] 
& \text{from Equation \eqref{little identity}}\\
&= \sum_j \epsilon_j [ E_j^\perp , \sum_i \epsilon_i \langle (\ad_{E_j^\perp})^* X,
E'_i\rangle E'_i]  &  \\
& = \sum_j \epsilon_j [E_j^\perp, (\ad_{E_j^\perp})^* X] & 
\text{because $X \in \frakn'$}\\ 
&= \sum_j \epsilon_j  (\ad_{E_j^\perp}) \circ  (\ad_{E_j^\perp})^* (X) 
= A. & \end{alignat*}

Next we prove $D=0$ by proving that each term in the sum vanishes.
We may assume $X \in \n_{\alpha}$ with $\alpha(Z) > 0$.  Take any $E'_i \in \n_{\beta}$ with $\beta(Z) > 0$. Then  $(\ad_{E'_i}) X \in \n_{\alpha + \beta}$. Hence $(\ad_{E'_i})^* (\ad_{E'_i}) X \in
\n_{\alpha} \subseteq \n'$, which implies its projection onto $\n_0$ vanishes: $\pi_{\n_0}((\ad_{E'_i})^* (\ad_{E'_i}) X) =0$. This shows that $D=0$.

Having shown that $A=B$ and $D=0$,  we see 
\[ \Ric^{\n} (X) - \Ric^{\n'} (X)  = \tfrac 14 A + \tfrac 14  B - \tfrac 12 C -\tfrac 12 D  = \tfrac 12 (A -  C).
\]
Substituting for $A$ and $C$ yields  the desired expression for $\Ric^{\n} (X) - \Ric^{\n'} (X) $.

We note that $\tfrac 12 (A -  C)$ is precisely the left-hand side of the Jacobi Star Condition in Equation \eqref{eq:jacobi-star}, while $\ad_{H-H'}(X)$ is precisely the right-hand side of the Jacobi Star Condition. Thus when the Jacobi Star Condition holds, we have
\[ \Ric^{\n} (X) - \Ric^{\n'} (X)  =  \ad_{H-H'}(X). 
\]
\end{proof}

\end{document}
